\tdplotsetmaincoords{70}{140}
\begin{tikzpicture}[tdplot_main_coords,scale=2.5]
    \clip(0,-1.6,-1) rectangle (0,1.6,1.6);

    %% Définition des différents styles
    \tikzstyle{init} = [black]			% base initiale
    \tikzstyle{prec} = [blue]			% 1ere base intermédiaire
    \tikzstyle{nuta} = [red]			% 2eme base initiale
    \tikzstyle{rotp} = [green]			% base finale
    \tikzstyle{base} = [thick,-stealth]	% Tracé des bases
    \tikzstyle{angle} = [thick,-latex]	% Tracé des arcs pour les angles
    \tikzstyle{circle} = [thin,dashed]	% Tracé des cercles
    \tikzstyle{dummy} = [draw=none,black]

    %% Paramètres géométriques
    \def\ephi{30}	% Angle de precession dessiné
    \def\etheta{45}	% Angle de nutation dessiné
    \def\rang{0.6}	% Rayon utilisé pour tracer les angles

    %% Tracé
    % Repère initial
    \coordinate (O) at (0,0,0);
    \node[anchor=south west] at (O) {$O$};
    \draw[base,init] (O) -- (1,0,0) node[anchor=north east]{$\overrightarrow{x}$};
    \draw[base,init] (O) -- (0,1,0) node[anchor=north west]{$\overrightarrow{y}$};
    \draw[base,init] (O) -- (0,0,1) node[anchor=south]{$\overrightarrow{z}$};

    % Précession
    \tdplotsetrotatedcoords{\ephi}{0}{0}
    \draw[tdplot_rotated_coords,angle,prec,dotted] (O) --(1,0,0) node[anchor=north east]{};
    \tdplotdrawarc[tdplot_rotated_coords,circle,prec]{(0,0,0)}{1}{0}{360}{}{}
    \tdplotdrawarc[tdplot_rotated_coords,angle,prec]{(0,0,0)}{\rang}{-\ephi}{0}{anchor=north east,prec}{$\phi$}

    % Nutation
    \tdplotsetrotatedcoords{\ephi}{90-\etheta}{0}
    \draw[tdplot_rotated_coords,base,nuta] (O) --(0,0,1) node[anchor=south east,init]{$P$};
    \draw[tdplot_rotated_coords,dummy] (O) -- (0,0,0.5) node[anchor=south,nuta]{$R$};
    \tdplotsetrotatedthetaplanecoords{0}
    \tdplotdrawarc[tdplot_rotated_coords,circle,nuta]{(0,0,0)}{1}{0}{360}{}{}
    \tdplotdrawarc[tdplot_rotated_coords,angle,nuta]{(0,0,0)}{\rang}{90-\etheta}{0}{anchor=west,nuta}{$\theta$}
\end{tikzpicture}